\section{Rank-one amplituhedral geometries and residue compatibility}
\label{sec:k-poligono-amplituhedral}

The dodecaphasic register allows positive data, Grassmannian domain,
kinematic image, covering, degree on cells, and canonical form to be
assembled explicitly. The first result concerns the two-dimensional
amplituhedron with parameters \((n,k,m)=(13,1,2)\); the construction
then extends to rank-one cyclic polytopes of higher dimension.
This is a subsequent geometric realization of the register produced
by APP--TRIT--TPK. The stated additional operation is lifting an ordered
rational partition by powers. This choice fixes a particular realization;
subsequent physical correspondences retain their own readers and conditions.

The provenance of the register preserves the composition
\[
 x_{\mathrm{term}}\longmapsto\mathcal R_{12}(x_{\mathrm{term}})
 \longmapsto(b^{90},b^{120},Q_{\mathrm{TPK}})
 \longmapsto U_{\mathrm{sgn}}\longmapsto K.
\]
APP preserves both sheets and their remainders and quotients; TRIT
orients the regime of each transition; TPK transports and updates
the state with carry, boundary, route, and memory. The representation
\(K\) is read in the same system of enriched states that jointly
generates the continuum. Its signed extraction and Hadamard inversion
produce
\begin{equation}
 K=(234,543,140,729,659,824,621,58,914,794,146,601),
 \qquad S_K=\sum_{j=1}^{12}K_j=6263.
 \label{eq:polygon-K}
\end{equation}
The geometric proof uses the positivity of these twelve coordinates,
their order, and their total. Its coefficients come from that extraction,
not from retrospective selection of a physical value.

\subsection{Rational partition and positive lift}

Define thirteen rational partition points:
\begin{equation}
 t_0=0,\qquad t_i=\frac{K_1+\cdots+K_i}{S_K}\quad(1\leq i\leq12),
 \qquad Z_i=(1,t_i,t_i^2)^\mathsf T,\qquad Z=(Z_0\ \cdots\ Z_{12}).
 \label{eq:polygon-Z}
\end{equation}
Their numerators, with common denominator \(6263\), are
\[
 (0,234,777,917,1646,2305,3129,3750,3808,4722,5516,5662,6263).
\]
Thus \(0=t_0<t_1<\cdots<t_{12}=1\). The labeled partition and
the total preserve the register through the exact inverse
\begin{equation}
 K_i=S_K(t_i-t_{i-1}).
 \label{eq:polygon-recover-K}
\end{equation}
The partition preserves the proportions; retaining \(S_K\) also
preserves the integer scale.

\begin{proposition}[Strictly positive kinematic data]
\label{prop:polygon-vandermonde}
The matrix \(Z\) has rank three, and all its ordered minors of size
three are positive. For \(0\leq a<b<c\leq12\),
\begin{equation}
 \langle abc\rangle:=\det(Z_a,Z_b,Z_c)
 =(t_b-t_a)(t_c-t_a)(t_c-t_b)>0.
 \label{eq:polygon-minors}
\end{equation}
\end{proposition}
\begin{proof}
Subtract the first column from the last two and expand along the
first row. Each difference \(t_j^2-t_a^2\) factors as
\((t_j-t_a)(t_j+t_a)\), giving the stated product. The strictly
increasing partial sums prove positivity. Any of these minors
establishes rank three.
\end{proof}

\subsection{Image, fibers, and degree-one cells}

The Grassmannian \(\operatorname{Gr}_{\geq0}(1,13)\) is the
projective space of nonnegative nonzero vectors
\(c=(c_0,\ldots,c_{12})\), modulo positive scaling.
The chart \(\sum c_j=1\) identifies it with the closed
twelve-dimensional simplex. Define
\begin{equation}
 \Phi_Z:\operatorname{Gr}_{\geq0}(1,13)\longrightarrow\mathbb P^2,
 \qquad[c]\longmapsto\left[\sum_{j=0}^{12}c_jZ_j\right].
 \label{eq:polygon-map}
\end{equation}
In the chart \(Y_0=1\), let \(v_j=(t_j,t_j^2)\) and
\(P_K=\operatorname{conv}(v_0,\ldots,v_{12})\).

\begin{theorem}[Polygonal amplituhedron of the register]
\label{thm:polygon-amplituhedron}
The closed image of \(\Phi_Z\) is exactly the convex polygon
\(P_K=\mathcal A_{13,1,2}(Z)\), of dimension two with thirteen
vertices. The image of the strictly positive Grassmannian is the
interior of \(P_K\). The eleven cells
\begin{equation}
 \mathcal S_i=\{[c]:c_0,c_i,c_{i+1}>0,\
 c_j=0\text{ if }j\notin\{0,i,i+1\}\},\qquad 1\leq i\leq11,
 \label{eq:polygon-cells}
\end{equation}
map bijectively, with degree one and positive orientation, onto the
interiors of \(T_i=[v_0,v_i,v_{i+1}]\). Their closures cover \(P_K\),
and their interiors are pairwise disjoint.
\end{theorem}
\begin{proof}
Under the normalization \(\sum c_j=1\), the map is
\(c\mapsto\sum c_jv_j\), whose image is the convex hull.
The minors in \eqref{eq:polygon-minors} show that every ordered
triple has positive orientation. Each consecutive side leaves all
other vertices strictly to its left; so does the final side from
\(v_{12}\) to \(v_0\). All points are therefore extreme vertices
of a convex polygon.

Every combination with \(c_j>0\) is interior, since for any
supporting line some vertex lies strictly in the interior half-plane.
Conversely, let \(y\) be interior and \(b=13^{-1}\sum v_j\).
For sufficiently small \(\varepsilon>0\),
\(y'=(y-\varepsilon b)/(1-\varepsilon)\) remains in the polygon.
Writing \(y'=\sum a_jv_j\), \(a_j\geq0\), \(\sum a_j=1\), gives
\(y=\sum((1-\varepsilon)a_j+\varepsilon/13)v_j\), with all
coefficients positive.

The diagonals from \(v_0\) divide the polygon into the eleven
triangles. On each triple support, the barycentric coordinates are
unique because \(\langle0,i,i+1\rangle>0\); their explicit inverse
appears in \eqref{eq:polygon-barycentric}. Each support defines a
rank-one positroid cell with coordinates \([1:a:b]\), \(a,b>0\),
at those positions. A directed planar network with one source and
three paths of weights \(1,a,b\) to the three destinations realizes
these coordinates by boundary measurement. The positive determinant
proves orientation and degree one.
\end{proof}

The generic fiber of the full map has dimension ten:
\(\sum c_j=1\), \(\sum c_jt_j=x\), and \(\sum c_jt_j^2=y\)
are three independent equations in thirteen coordinates.
For interior \(y\), the fiber contains a strictly positive point,
so its intersection with the simplex has dimension \(13-3=10\).
Degree one concerns the triangular cells; the two-dimensional form
is obtained through these cells.

\subsection{Canonical form and oriented residues}

For \(Y=(1,x,y)^\mathsf T\), define
\begin{equation}
 \ell_{ab}(Y)=\det(Z_a,Z_b,Y)
 =(t_b-t_a)\{y-(t_a+t_b)x+t_at_b\}.
 \label{eq:polygon-lines}
\end{equation}
The expression also holds when \(b<a\). For \(a<b<c\), let
\(D_{abc}=\langle abc\rangle\) and
\begin{equation}
 \lambda_a=\frac{\ell_{bc}}{D_{abc}},\qquad
 \lambda_b=\frac{\ell_{ca}}{D_{abc}},\qquad
 \lambda_c=\frac{\ell_{ab}}{D_{abc}},\qquad
 \lambda_a+\lambda_b+\lambda_c=1.
 \label{eq:polygon-barycentric}
\end{equation}
The affine orientation is \(dx\wedge dy\). Fix the boundary convention
\(\operatorname{Res}^{\partial}(\eta\wedge du/u)=\eta|_{u=0}\):
the normal differential occupies the last position. In dimension two,
the Poincaré residue placing \(du/u\) first has the opposite sign.
The following cancellations uniformly use the stated convention.

\begin{theorem}[Rational form and cancellation of diagonals]
\label{thm:polygon-form}
The canonical form of the oriented triangle \([v_a,v_b,v_c]\) is
\begin{equation}
 \Omega_{abc}=
 \frac{D_{abc}^2\,dx\wedge dy}
 {\ell_{ab}(Y)\ell_{bc}(Y)\ell_{ca}(Y)}.
 \label{eq:polygon-triangle-form}
\end{equation}
The canonical form of the polygon is
\begin{equation}
 \boxed{\displaystyle
 \Omega_{P_K}(x,y)=
 \sum_{i=1}^{11}
 \frac{t_it_{i+1}(t_{i+1}-t_i)\,dx\wedge dy}
 {(y-t_ix)\{y-(t_i+t_{i+1})x+t_it_{i+1}\}(t_{i+1}x-y)}.}
 \label{eq:polygon-full-form}
\end{equation}
The poles of interior diagonals cancel. The only polar divisors are
the thirteen lines of the exterior sides. Parametrizing each side
from its initial to its final vertex by \(s\in[0,1]\), its boundary
residue is \(ds/[s(1-s)]\).
\end{theorem}
\begin{proof}
The form of the barycentric simplex is
\(d\lambda_b\wedge d\lambda_c/(\lambda_a\lambda_b\lambda_c)\).
The Jacobian satisfies
\(dx\wedge dy=D_{abc}\,d\lambda_b\wedge d\lambda_c\).
Substituting \eqref{eq:polygon-barycentric} gives
\eqref{eq:polygon-triangle-form}. For \(a=0\), canceling the three
scalar factors of the lines gives each summand of
\eqref{eq:polygon-full-form}.

On side \(ab\), set \(\lambda_c=u\),
\(\lambda_b=(1-u)s\), \(\lambda_a=(1-u)(1-s)\). Then
\[
 \Omega_{abc}=\frac{ds}{s(1-s)}\wedge\frac{du}{u(1-u)}.
\]
The residue is the form of the oriented interval. Two adjacent
triangles induce opposite orientations on their diagonal and produce
opposite residues. Since all poles are simple, a vanishing residue
removes the diagonal pole as a divisor. Each exterior side appears
once. The homogeneous formula has degree zero as a projective form
and has no other polar divisors. If two rational projective forms
had the same residues and no other poles, their difference would be
a holomorphic two-form on \(\mathbb P^2\), necessarily zero.
This proves characterization and uniqueness.
\end{proof}

We use the conventional notion of positive geometry through logarithmic
poles and canonical boundary residues, developed by
N.~Arkani-Hamed, Y.~Bai, and T.~Lam, \emph{Positive Geometries and Canonical
Forms}, JHEP 11 (2017), 039,
\url{https://arxiv.org/abs/1703.04541}, Sections 2, 3, 5.2, 6.1, and 7.2.
The selection of the coefficients of the present polygon precedes this
recognition and comes from \(K\).

\subsection{Carry refinement on the fixed geometry}

In a positively oriented triangle \(T=[A,B,C]\), write
\begin{equation}
 Y(u,r)=(1-u)A+u\{(1-r)B+rC\},\qquad 0<u,r<1.
 \label{eq:polygon-blowdown}
\end{equation}
Vertex \(A\) corresponds to \(u=0\), and the opposite side to
\(u=1\). The form becomes
\begin{equation}
 \Omega_T=\frac{du}{u(1-u)}\wedge\frac{dr}{r(1-r)}.
 \label{eq:polygon-product-form}
\end{equation}
Let \(B_0\geq2\) be a stated integer base, distinct from vertex
\(B\). The coordinate \(r\) admits the remainder-and-carry update
\begin{equation}
 c=\lfloor B_0r\rfloor,\qquad r'=B_0r-c,\qquad
 r=\frac{c+r'}{B_0}.
 \label{eq:polygon-carry}
\end{equation}
The endpoints retain their boundary marks. The digit \(c\) remains
as memory and allows the update to be inverted. The charts of the
two affine transports should be specified. The subcell zoom acts
on \(r\in[c/B_0,(c+1)/B_0]\) through \(r'=B_0r-c\).
By contrast, transport of the lifted pair
\(L_c(x,y)=(B_0x-c,B_0y+c)\), with \(x+y=M\), changes the
total mass to \(B_0M\) and expresses the coordinate
\[
 \frac{B_0x-c}{B_0M}=\frac{x}{M}-\frac{c}{B_0M}.
\]
These maps have different domains and normalizations. Both preserve
the carry in the enriched state and admit recovery with their scale
data. The following proof uses the subcell zoom and the geometric
partition it identifies.

\begin{theorem}[Exact invariance under subdivision with memory]
\label{thm:polygon-carry}
Let \(s_j=j/B_0\) and \(W_j=(1-s_j)B+s_jC\).
The triangles \(T_j=[A,W_j,W_{j+1}]\), \(0\leq j<B_0\),
subdivide \(T\) and satisfy
\begin{equation}
 \sum_{j=0}^{B_0-1}\Omega_{T_j}=\Omega_T.
 \label{eq:polygon-carry-additivity}
\end{equation}
The equality remains valid at every finite depth and for any ordered
rational partition of the side. Subdividing the eleven triangles
of \(P_K\) leaves \(\Omega_{P_K}\) exactly invariant.
\end{theorem}
\begin{proof}
The regions \(r\in[s_j,s_{j+1}]\) cover \(T\) with disjoint
interiors. In the common coordinates,
\[
 \Omega_{T_j}=\frac{du}{u(1-u)}\wedge
 \frac{(s_{j+1}-s_j)\,dr}{(r-s_j)(s_{j+1}-r)}.
\]
The rational identity
\[
 \frac{s_{j+1}-s_j}{(r-s_j)(s_{j+1}-r)}
 =\frac1{r-s_j}+\frac1{s_{j+1}-r}
\]
telescopes: each interior term cancels the one from the adjacent
subinterval. This leaves \(1/r+1/(1-r)=1/[r(1-r)]\).
This proves equality of rational forms; iteration gives every finite
depth. The proof uses only \(s_0=0<s_1<\cdots<s_N=1\), so it
covers nonuniform partitions. The change
\(r'=(r-s_j)/(s_{j+1}-s_j)\) locally returns the unit form.
The carry word identifies the subcell, and its inverse recovers
the previous coordinate. The new interior poles cancel through the
same telescoping identity.
\end{proof}

The preceding sum assembles the complete geometric partition. When a
survival reader selects only some histories, its image is the union
of the selected subcells and its form is the corresponding sum;
it agrees with the form of the full polygon precisely when that union
covers it with oriented multiplicity one. The covering criterion for
a specific selection of histories is therefore expressed by a
checkable condition on its cell chain.

This refinement preserves the fixed polygon. Inserting a new point
\((t,t^2)\) on the parabola between consecutive parameters
\(a<t<b\) produces another geometry: the chord height minus
\(t^2\) is \((t-a)(b-t)>0\), and the new point lies outside
the previous polygon. The proved invariance concerns rectilinear
subdivision of cells, whose exterior boundary remains fixed.

\subsection{Triangulation changes and positive reparametrization}

For four ordered vertices \(A,B,C,D\), we have
\begin{equation}
 \Omega_{ABC}+\Omega_{ACD}=\Omega_{ABD}+\Omega_{BCD}.
 \label{eq:polygon-flip}
\end{equation}
Both sides have the same exterior residues and cancel their interior
diagonal; their equality follows from the proved uniqueness.
It is also checked by multiplying through by their linear denominators.
Every change between triangulations of a convex polygon decomposes
into these diagonal changes, so the form is independent of triangulation.
Identification with cluster coordinates must additionally preserve
the variety and its change relations.

If \(D=\operatorname{diag}(d_0,\ldots,d_{12})\), \(d_j>0\), then
\begin{equation}
 \mathcal A_{13,1,2}(ZD)=\mathcal A_{13,1,2}(Z).
 \label{eq:polygon-rescale}
\end{equation}
The automorphism \([c]\mapsto[Dc]\) reparametrizes the kinematic map.
In the affine chart, \(c_j\) changes to
\(c_jd_j/\sum_\ell c_\ell d_\ell\), which again ranges over the
simplex. The image and its canonical form remain unchanged.
Displacement of a parametrized history and change of the entire
geometric set are thus distinguishable operations.

\subsection{Rank-one extension to higher dimensions}

For \(2\leq m\leq12\), define
\begin{equation}
 Z_i^{(m)}=(1,t_i,t_i^2,\ldots,t_i^m)^\mathsf T,\qquad
 P_K^{(m)}=\operatorname{conv}
 \{(t_i,t_i^2,\ldots,t_i^m):0\leq i\leq12\}.
 \label{eq:polygon-higher-lift}
\end{equation}
The ordered maximal minors are the Vandermonde products
\(\prod_{a<b}(t_{i_b}-t_{i_a})>0\). The matrix has rank \(m+1\).
The image of \(\operatorname{Gr}_{\geq0}(1,13)\) is the cyclic
polytope \(P_K^{(m)}=\mathcal A_{13,1,m}(Z^{(m)})\), of dimension
\(m\); its generic fiber has dimension \(12-m\). Here we use the
mathematical extension of the amplituhedral definition to any \(m\);
customary applications to amplitudes impose their own values of \(m\).

The triangulation is determined by the fixed label order. Choose the
least vertex of a polytope, recursively triangulate all facets not
containing it, and take cones from that vertex. Induction begins with
points. A line from the initial vertex to an interior point meets
the opposite boundary in a facet; this proves covering. Uniqueness
of the intersection point and compatibility of the restrictions
induced by the same order prove that interiors are disjoint and
intersections are common faces. Dimension decreases at each step,
so the procedure terminates. This is a pulling triangulation
determined by the labels.

For an oriented simplex \(T=[p_0,\ldots,p_m]\) of this triangulation,
let
\[
 A_T=(p_1-p_0\ \cdots\ p_m-p_0),\qquad
 (\lambda_1,\ldots,\lambda_m)^\mathsf T=A_T^{-1}(x-p_0),
 \qquad\lambda_0=1-\sum_{j=1}^m\lambda_j.
\]
Orient its vertices so that \(\det A_T>0\).
Its canonical form is
\begin{equation}
 \Omega_T(x)=
 \frac{dx_1\wedge\cdots\wedge dx_m}
 {\det(A_T)\prod_{j=0}^m\lambda_j(x)},\qquad
 \Omega_{P_K^{(m)}}=\sum_{T}\Omega_T.
 \label{eq:polygon-higher-form}
\end{equation}
The barycentric change proves the formula and that the kinematic
restriction to the support of \(T\) has degree one. The residues
are the face forms with signs induced by orientation. On an interior
face, the two orientations are opposite and the total residue vanishes.
On an exterior face, the sum of the forms in that face's triangulation
remains. Induction on dimension proves that this sum is its canonical
form. Repeating the residue along a flag of faces reaches the vertices,
where the normalization is \(\pm1\). This gives residue correspondence
in every codimension for this rank-one collection.

The same proof covers any compatible simplicial subdivision of the
fixed polytope: interior faces cancel and exterior faces preserve the
residues. Nonexistence of top-degree holomorphic forms on
\(\mathbb P^m\) ensures uniqueness of the rational form.
In particular, refinement stability, covering, and residues are
settled in this collection of dimensions \(2,\ldots,12\).
Grassmannian degrees \(k>1\) retain their specific maps and
proof obligations.
